\section{From Relative Signals to the Next Song}
\label{sec:method}

\subsection{Design principle}

Vibe Shuffle is a client-only React/Vite application. A camera, an optional
Bluetooth Heart Rate Service device, and Spotify playback meet inside one
browser session. There is no study database or recommendation server. The
design keeps adaptation between songs rather than changing playback
continuously: the system listens to the latest window, summarizes it, and then
chooses the next track. This makes the decision legible and prevents every
short-lived fluctuation from becoming a command.

Figure~\ref{fig:trajectory} shows the complete idea. Two personal references
anchor the incoming signals. Confirmed playback produces a path rather than a
single instantaneous label. The path is averaged, mapped to a region, and used
to select a nearby song. A participant then rates the track, closing the
interaction loop.

\begin{figure*}[t]
  \centering
  \includegraphics[width=\textwidth]{fig2_relative_trajectory.pdf}
  \caption{\textbf{Relative sensing, trajectory, and selection in Vibe
  Shuffle.} (a) Camera and cardiac inputs are expressed as signed changes from
  personal references. (b) A curved Valence--Arousal trajectory is summarized
  by its arithmetic mean (gold diamond); points grow and darken from the open
  start marker to the filled endpoint. (c) Within the matching region,
  \emph{As It Was} is nearest but
  excluded as played (cross), so \emph{CUFF IT} becomes the selected eligible
  track (filled diamond); open circles are alternatives.
  The relative changes and trajectory are illustrative; catalog candidates,
  distances, and the selection result are repository-derived.
  This is a relative heuristic, not an emotion diagnosis; coherence is logged
  separately and does not affect selection.}
  \Description{Three equal square panels show camera and cardiac changes from
  personal references, a curved illustrative relative valence-arousal
  trajectory with its mean, and a ranked list of five catalog candidates in
  which As It Was is marked played and CUFF IT is selected.}
  \label{fig:trajectory}
\end{figure*}

\subsection{Relative Valence from camera cues}

The browser runs MediaPipe Face Landmarker locally~\cite{mediapipe2019}. It
returns facial blendshape activations such as smile, cheek, brow, eye, jaw,
frown, and mouth movement. Vibe Shuffle combines a selected subset into
transparent positive, negative, and tension-related scores. These labels belong
to Vibe Shuffle's rule layer; MediaPipe itself does not output the project's
emotion-region labels.

A 14-second calibration captures a short facial reference, and a slow running
baseline continues to adapt. Later observations are debiased against this
reference and smoothed across frames. Positive facial evidence moves Valence to
the right; negative or tension-related evidence moves it to the left. If no face
is detected, Valence returns to the neutral center. The current heuristic also
lets visible movement increase positive Valence evidence and upward Arousal
evidence. Movement may instead reflect dancing, repositioning, lighting change,
or camera noise, so this rule is not a validated engagement or affect measure.

\subsection{Relative Arousal from cardiac timing}

The optional wearable provides device-produced beats per minute and, when
available, RR intervals---the elapsed time between successive beats. The
application does not receive a raw electrocardiogram and cannot inspect the
device's peak-detection process. Implausible and abruptly changing RR intervals
are rejected before short-term variability is calculated.

The primary variability feature is the root mean square of successive
differences (RMSSD). A 120-second reference stores the listener's ordinary
median heart rate and typical RMSSD together with robust estimates of their
variation. During listening, higher-than-reference heart rate and
lower-than-reference RMSSD provide upward Arousal evidence. Heart rate receives
the larger influence because RMSSD estimated from a short window is unstable;
when the two disagree, the RMSSD contribution is strongly reduced. At least 20
accepted RR intervals are required before the full estimate is treated as
usable. Controlled resting data show high agreement between 120-second and
longer RMSSD recordings~\cite{munoz2015}; this supports the calibration duration,
not short moving-wearable windows or Arousal validity.

The implementation also exports SDNN and a custom normalized
frequency-domain spectral-peakedness score. The CSV calls this the physiology
coherence field. It is an experimental diagnostic, not the coherence measure
analyzed by Balaji et al., and it does not determine the selected region. Prior
HRV-biofeedback research reports group-level associations between coherence
patterns and end-of-session self-reported emotion
categories~\cite{balaji2025}. That result motivates careful study of cardiac
timing, but it neither validates the prototype's proxy nor shows that a short
personal window can identify an individual's emotion. Vibe Shuffle accordingly
uses baseline-relative activation evidence rather than translating coherence
frequencies into labels.

\subsection{A trajectory rather than a snapshot}

Whenever the fused live coordinate changes during confirmed playback, Vibe
Shuffle appends a point to the current trial buffer. Facial cues primarily
control Valence. Usable cardiac evidence supplies the Arousal base, with visible
movement able to add upward Arousal. If cardiac evidence is unavailable, camera
movement provides an upward-only fallback. If neither channel is usable, the
coordinate returns to $(0.5,0.5)$.

The resulting sequence makes the listener's path visible conceptually:
\[
  \mathcal{T}=\{(V_1,A_1),\ldots,(V_m,A_m)\}.
\]
At rating time, the application uses the arithmetic mean
\[
  (\bar V,\bar A)=\frac{1}{m}\sum_{i=1}^{m}(V_i,A_i)
\]
as the target for the next adaptive choice. The points are event-driven rather
than sampled at a fixed rate, so this is not a time-weighted physiological
average. Sample count and actual listening duration are exported to make that
approximation visible.

\subsection{Four regions and a fixed music space}

Splitting both axes at 0.5 produces four recommendation regions:
Energetic (higher Valence and Arousal), Calm (higher Valence and lower
Arousal), Tense (lower Valence and higher Arousal), and Melancholic (lower on
both axes). The frozen catalog contains 100 Spotify tracks, with 25 tracks in
each corresponding Valence--Energy region. Track coordinates originate from
historical Spotify features in a public dataset~\cite{spotifydataset}. Energy
is a practical catalog-side counterpart of Arousal, not proof that an audio
feature and a physiological construct are identical.

Both Random and Vibe exclude the current and all previously played tracks.
Random ranks the remaining tracks with a session-seeded pseudorandom score.
Vibe first retains available tracks in the target region and then selects the
smallest Euclidean distance:
\[
 s^*=\arg\min_{s\in Q(\bar V,\bar A)}
 \sqrt{(V_s-\bar V)^2+(E_s-\bar A)^2}.
\]
The seed resolves only an exact distance tie. If a region has no remaining
track, the policy falls back to the complete unplayed pool. The export records
the target coordinate, signal source, distance, and region match. It does not
record the random seed, the personal facial-reference values, or individual
trajectory points. The first track is also chosen from the live setup estimate
before the 14-second facial reference is finalized; later choices use the
preceding listening window.

\subsection{Browser implementation and privacy boundary}

Spotify authentication uses Authorization Code with PKCE, and the Web Playback
SDK provides playback for an allowlisted Premium account. Raw camera frames,
facial landmarks, and Bluetooth notifications remain in browser memory. The
locally downloaded CSV contains derived signal summaries, song metadata,
selection fields, ratings, quality flags, pseudonymous identifiers, and
timestamps. It is downloaded automatically at protocol completion and can also
be downloaded manually. It contains no image, landmark set, raw waveform, or
audio. Because it retains only trajectory summaries, it supports review of the
final adaptive target but not replay of every sensing step.

This is a bounded privacy property, not a claim of offline operation. Spotify
authentication and playback, MediaPipe asset delivery, and GitHub Pages create
ordinary external traffic. Exported participant numbers and timestamps are
pseudonymous rather than anonymous and still require appropriate research-data
handling.
